\documentclass[12pt,a4paper]{report}

% ── Packages ──────────────────────────────────────────────────────────────────
\usepackage[a4paper, left=1.5in, right=1in, top=1in, bottom=1in]{geometry}
\usepackage{times}
\usepackage{setspace}
\usepackage{graphicx}
\usepackage{booktabs}
\usepackage{array}
\usepackage{longtable}
\usepackage{tabularx}
\usepackage{multirow}
\usepackage{fancyhdr}
\usepackage{titlesec}
\usepackage{tocloft}
\usepackage{hyperref}
\usepackage{listings}
\usepackage{xcolor}
\usepackage{float}
\usepackage{enumitem}
\usepackage{indentfirst}
\usepackage{caption}
\usepackage{amsmath}
\usepackage{amssymb}
\usepackage[T1]{fontenc}
\usepackage{microtype}

% ── Hyperref setup ────────────────────────────────────────────────────────────
\hypersetup{
    colorlinks=true,
    linkcolor=black,
    filecolor=black,
    urlcolor=black,
    citecolor=black,
}

% ── Line spacing ──────────────────────────────────────────────────────────────
\doublespacing

% ── Chapter title formatting (match Doc.pdf) ──────────────────────────────────
\titleformat{\chapter}[block]
  {\normalfont\Large\bfseries\centering}
  {Chapter \thechapter}{0pt}{\\\Large\bfseries\centering}
\titlespacing*{\chapter}{0pt}{30pt}{20pt}

\titleformat{\section}[hang]
  {\normalfont\normalsize\bfseries}
  {\thesection}{1em}{}
\titleformat{\subsection}[hang]
  {\normalfont\normalsize\bfseries}
  {\thesubsection}{1em}{}
\titleformat{\subsubsection}[hang]
  {\normalfont\normalsize\bfseries}
  {\thesubsubsection}{1em}{}

% ── Header/Footer ─────────────────────────────────────────────────────────────
\pagestyle{fancy}
\fancyhf{}
\fancyfoot[C]{\thepage}
\renewcommand{\headrulewidth}{0pt}

% ── Code listing style ────────────────────────────────────────────────────────
\lstset{
  basicstyle=\small\ttfamily,
  breaklines=true,
  frame=single,
  columns=flexible,
}

% ── TOC depth ─────────────────────────────────────────────────────────────────
\setcounter{tocdepth}{3}
\setcounter{secnumdepth}{3}

% ══════════════════════════════════════════════════════════════════════════════
\begin{document}

% ── PAGE i: COVER PAGE ────────────────────────────────────────────────────────
\begin{titlepage}
\centering
\vspace*{1cm}
{\Large\textbf{OPINION MINER}\\[6pt]
A MINI PROJECT REPORT}\\[1cm]
{\normalsize Submitted by}\\[6pt]
{\large\textbf{Bandaru Bhanu Prasad (B210054)\\
Eslavath Ashok (B210474)\\
Thagilla Rakesh (B211571)}}\\[1cm]
{\normalsize\textbf{Of}}\\[6pt]
{\large\textbf{Bachelor of Technology}}\\[6pt]
{\normalsize Under the guidance of}\\[4pt]
{\large\textbf{Miss.\ Riharika , Asst.\ Prof}}\\[1.5cm]

% University logo placeholder
\begin{figure}[H]
\centering
\fbox{\parbox{3cm}{\centering\vspace{1.5cm}University\\Logo\vspace{1.5cm}}}
\end{figure}
% Replace above with: \includegraphics[width=3cm]{rgukt_logo.png}

\vspace{0.5cm}
{\small\textbf{RGUKT BASAR}}\\[6pt]
{\large\textbf{DEPARTMENT OF COMPUTER SCIENCE AND ENGINEERING\\
RAJIV GANDHI UNIVERSITY OF KNOWLEDGE TECHNOLOGIES\\
BASAR, NIRMAL (DIST.), TELANGANA -- 504107\\
MARCH 2026}}

\thispagestyle{fancy}
\fancyfoot[C]{\roman{page}}
\setcounter{page}{1}
\end{titlepage}

% ── PAGE ii: INNER TITLE PAGE ─────────────────────────────────────────────────
\newpage
\pagenumbering{roman}
\setcounter{page}{2}
\thispagestyle{fancy}
\fancyfoot[C]{\roman{page}}

\begin{center}
\vspace*{1.5cm}
{\Large\textbf{OPINION MINER}}\\[1cm]
{\normalsize\textit{Project Report submitted to\\
Rajiv Gandhi University of Knowledge Technologies, Basar\\
for the partial fulfillment of the requirements\\
for the award of the degree of}}\\[0.8cm]
{\large\textbf{Bachelor of Technology}}\\[4pt]
{\normalsize in}\\[4pt]
{\large\textbf{Computer Science \& Engineering}}\\[0.8cm]
{\normalsize\textit{by}}\\[0.5cm]
{\large\textbf{\textit{Bandaru Bhanu Prasad (B210054)\\
Eslavath Ashok (B210474)\\
Thagilla Rakesh (B211571)}}}\\[0.8cm]
{\large\textbf{Under the Guidance of\\
\textit{Miss.\ Riharika , Asst.\ Prof}}}\\[1.5cm]

\begin{figure}[H]
\centering
\fbox{\parbox{3cm}{\centering\vspace{1.5cm}University\\Logo\vspace{1.5cm}}}
\end{figure}

\vspace{0.4cm}
{\small\textbf{RGUKT BASAR}}\\[4pt]
{\large\textbf{DEPARTMENT OF COMPUTER SCIENCE AND ENGINEERING\\
RAJIV GANDHI UNIVERSITY OF KNOWLEDGE TECHNOLOGIES, BASAR\\
MARCH 2026}}
\end{center}

% ── PAGE iii: CERTIFICATE ─────────────────────────────────────────────────────
\newpage
\setcounter{page}{3}
\thispagestyle{fancy}
\fancyfoot[C]{\roman{page}}

\begin{center}
\begin{figure}[H]
\centering
\fbox{\parbox{3cm}{\centering\vspace{1.5cm}University\\Logo\vspace{1.5cm}}}
\end{figure}
{\small\textbf{RGUKT BASAR\\
DEPARTMENT OF COMPUTER SCIENCE AND ENGINEERING\\
RAJIV GANDHI UNIVERSITY OF KNOWLEDGE TECHNOLOGIES, BASAR}}\\[1cm]
{\large\textbf{\underline{CERTIFICATE}}}\\[0.8cm]
\end{center}

\noindent This is to certify that the Project Report entitled \textbf{Opinion Miner -- A Smart System for Analyzing Online Product Reviews} submitted by \textbf{Bandaru Bhanu Prasad (B210054), Eslavath Ashok (B210474), Thagilla Rakesh (B211571)}, Department of Computer Science and Engineering, Rajiv Gandhi University Of Knowledge Technologies, Basar. For partial fulfillment of the requirements for the degree of Bachelor of Technology in Computer Science and Engineering; is a bonafide record of the work and investigations carried out by them under my supervision and guidance.

\vspace{2.5cm}
\noindent
\begin{minipage}[t]{0.45\textwidth}
\textbf{PROJECT SUPERVISOR}\\[1.2cm]
\textbf{Miss.\ Riharika}\\
Assistant Professor
\end{minipage}
\hfill
\begin{minipage}[t]{0.45\textwidth}
\textbf{HEAD OF DEPARTMENT}\\[1.2cm]
\textbf{Dr.\ B.\ Venkat Raman}\\
Assistant Professor
\end{minipage}

\vspace{2cm}
\begin{center}
\textbf{EXTERNAL EXAMINER}
\end{center}

% ── PAGE iv: DECLARATION ──────────────────────────────────────────────────────
\newpage
\setcounter{page}{4}
\thispagestyle{fancy}
\fancyfoot[C]{\roman{page}}

\begin{center}
\begin{figure}[H]
\centering
\fbox{\parbox{3cm}{\centering\vspace{1.5cm}University\\Logo\vspace{1.5cm}}}
\end{figure}
{\small\textbf{RGUKT\\
DEPARTMENT OF COMPUTER SCIENCE AND ENGINEERING\\
RAJIV GANDHI UNIVERSITY OF KNOWLEDGE TECHNOLOGIES, BASAR}}\\[0.8cm]
{\large\textbf{\underline{DECLARATION}}}\\[0.8cm]
\end{center}

\noindent We hereby declare that the work which is being presented in this project entitled, ``Opinion Miner -- A Smart System for Analyzing Online Product Reviews'' submitted to RAJIV GANDHI UNIVERSITY OF KNOWLEDGE TECHNOLOGIES, BASAR in the partial fulfillment of the requirements for the award of the degree of BACHELOR OF TECHNOLOGY in COMPUTER SCIENCE AND ENGINEERING, is an authentic record of our own work carried out under the supervision of ``Miss.\ Riharika'', Assistant Professor in Department of Computer Science And Engineering, RGUKT, Basar.

\medskip
\noindent The matter embodied in this project report has not been submitted by us for the award of any other degree.

\vspace{1.5cm}
\noindent Place: Basar

\noindent Date : 26-03-2026

\vspace{1.5cm}
\begin{flushright}
Bandaru Bhanu Prasad (B210054)\\
Eslavath Ashok (B210474)\\
Thagilla Rakesh (B211571)
\end{flushright}

% ── PAGE v: ACKNOWLEDGEMENT ───────────────────────────────────────────────────
\newpage
\setcounter{page}{5}
\thispagestyle{fancy}
\fancyfoot[C]{\roman{page}}

\begin{center}
{\large\textbf{\underline{ACKNOWLEDGEMENT}}}
\end{center}

\medskip
\noindent\hspace{1cm} We would like to express our sincere gratitude to our project advisor, Miss.\ Riharika, whose knowledge and guidance has motivated us to achieve our goals. She has consistently been a source of motivation, encouragement, and inspiration. The time we have spent working under her supervision has truly been a pleasure.

\medskip
\noindent\hspace{1cm} We thank HOD Dr.\ B.\ Venkat Raman for his effort and guidance and all senior faculty members of the CSE Department for their help during our course. Thanks to the programmers and non-teaching staff of the CSE Department for their continued support.

\medskip
\noindent\hspace{1cm} We would like to express our sincere gratitude to our Project Co-ordinator Mrs.\ B.\ Latha Reddy, whose knowledge and guidance has motivated us to achieve goals we never thought possible. She has consistently been a source of motivation, encouragement, and inspiration. The time we have spent working under her supervision has truly been a pleasure.

\medskip
\noindent\hspace{1cm} We thank our Vice-Chancellor, Prof.\ A.\ Govardhan, and the Management for providing excellent facilities to carry out our project work.

\medskip
\noindent\hspace{1cm} Finally, special thanks to our parents for their support and encouragement throughout our lives and this course. Thanks to all our friends and well-wishers for their constant support.

\vspace{1.5cm}
\begin{flushright}
Bandaru Bhanu Prasad (B210054)\\
Eslavath Ashok (B210474)\\
Thagilla Rakesh (B211571)
\end{flushright}

% ── PAGE vi: ABSTRACT ─────────────────────────────────────────────────────────
\newpage
\setcounter{page}{6}
\thispagestyle{fancy}
\fancyfoot[C]{\roman{page}}

\begin{center}
{\large\textbf{Abstract}}
\end{center}

\medskip
\noindent\hspace{1cm} In today's rapidly evolving digital landscape, e-commerce platforms such as Flipkart receive millions of product reviews daily. Understanding customer sentiment hidden within this massive volume of unstructured text data is crucial for both consumers making purchasing decisions and businesses strategizing product improvements. However, manual analysis of these reviews is not only impractical but also highly susceptible to bias, inconsistency, and the inability to process large-scale data in real time.

\medskip
\noindent\hspace{1cm} Opinion Miner is an intelligent, Flask-based web application designed to analyze sentiment in both raw user-entered text and Flipkart product reviews through automated scraping, text preprocessing, and multi-model sentiment analysis. The system employs the VADER (Valence Aware Dictionary and sEntiment Reasoner) model to compute positive, negative, neutral, and compound sentiment scores for individual and aggregated reviews. It integrates a Selenium-based web scraper capable of extracting real-time product reviews from Flipkart, circumventing dynamic page rendering barriers, thereby enabling live analysis.

\medskip
\noindent\hspace{1cm} A key innovation of the system is its Machine Learning-based Fake Review Detection module powered by a LinearSVC Support Vector Machine trained on 40,432 labelled reviews from the Cornell Deceptive Opinion Spam dataset, achieving a classification accuracy of 91.55\%. By eliminating unreliable reviews from the analysis pipeline, the system ensures that only genuine customer feedback influences the final sentiment outcome.

\medskip
\noindent\hspace{1cm} The application further enhances its analytical depth through Aspect-Based Sentiment Analysis (ABSA), which decomposes overall product sentiment into nine granular product dimensions: Quality, Delivery, Value, Performance, Battery, Camera, Display, Design, and Service. Each aspect is scored independently using sentence-level VADER analysis, enabling actionable product insights beyond a single aggregate score.

\medskip
\noindent\hspace{1cm} A Chrome Browser Extension enables direct sentiment analysis on any open Flipkart product page without navigating away. Visualization of results through bar charts and word clouds is powered by Matplotlib and WordCloud libraries. The system is containerised using Docker and deployed on Render.com for public accessibility.

\medskip
\noindent\hspace{1cm} The platform is developed using Python as the core programming language, with Flask as the backend web framework, Selenium and BeautifulSoup for scraping, scikit-learn for machine learning, HTML, CSS, and JavaScript for the front-end. The system successfully bridges the gap between raw, unstructured e-commerce feedback and actionable, trustworthy consumer intelligence.

% ── TABLE OF CONTENTS ─────────────────────────────────────────────────────────
\newpage
\setcounter{page}{7}
\thispagestyle{fancy}
\fancyfoot[C]{\roman{page}}
\renewcommand{\contentsname}{\centering Contents}
\tableofcontents

% ── LIST OF FIGURES ───────────────────────────────────────────────────────────
\newpage
\setcounter{page}{9}
\thispagestyle{fancy}
\fancyfoot[C]{\roman{page}}
\renewcommand{\listfigurename}{\centering List of Figures}
\listoffigures

% ── LIST OF ABBREVIATIONS ─────────────────────────────────────────────────────
\newpage
\setcounter{page}{10}
\thispagestyle{fancy}
\fancyfoot[C]{\roman{page}}

\begin{center}
{\large\textbf{List of Abbreviations and Symbols Used}}
\end{center}

\vspace{0.5cm}
\begin{longtable}{|p{3cm}|p{10cm}|}
\hline
\textbf{NLP}    & Natural Language Processing \\ \hline
\textbf{VADER}  & Valence Aware Dictionary and sEntiment Reasoner \\ \hline
\textbf{ML}     & Machine Learning \\ \hline
\textbf{SVM}    & Support Vector Machine \\ \hline
\textbf{TF-IDF} & Term Frequency-Inverse Document Frequency \\ \hline
\textbf{ABSA}   & Aspect-Based Sentiment Analysis \\ \hline
\textbf{API}    & Application Programming Interface \\ \hline
\textbf{UI}     & User Interface \\ \hline
\textbf{UX}     & User Experience \\ \hline
\textbf{HTML}   & HyperText Markup Language \\ \hline
\textbf{CSS}    & Cascading Style Sheets \\ \hline
\textbf{JS}     & JavaScript \\ \hline
\textbf{HTTP}   & HyperText Transfer Protocol \\ \hline
\textbf{URL}    & Uniform Resource Locator \\ \hline
\textbf{CSV}    & Comma Separated Values \\ \hline
\textbf{IDE}    & Integrated Development Environment \\ \hline
\textbf{DOM}    & Document Object Model \\ \hline
\textbf{BOW}    & Bag of Words \\ \hline
\textbf{BERT}   & Bidirectional Encoder Representations from Transformers \\ \hline
\textbf{RAM}    & Random Access Memory \\ \hline
\textbf{SSD}    & Solid State Drive \\ \hline
\textbf{HDD}    & Hard Disk Drive \\ \hline
\textbf{OS}     & Operating System \\ \hline
\textbf{PDF}    & Portable Document Format \\ \hline
\textbf{MVC}    & Model-View-Controller \\ \hline
\textbf{RNW}    & React Native Web \\ \hline
\textbf{WSGI}   & Web Server Gateway Interface \\ \hline
\textbf{JSON}   & JavaScript Object Notation \\ \hline
\textbf{CG}     & Computer Generated (Fake Review Label) \\ \hline
\textbf{OR}     & Original Review (Genuine Review Label) \\ \hline
\end{longtable}

% ══════════════════════════════════════════════════════════════════════════════
%  CHAPTER 1: INTRODUCTION
% ══════════════════════════════════════════════════════════════════════════════
\newpage
\pagenumbering{arabic}
\setcounter{page}{1}
\fancyfoot[C]{\thepage}

\chapter{Introduction}

\section{INTRODUCTION}

\noindent\hspace{1cm} In the modern digital era, e-commerce platforms have transformed the way consumers shop, evaluate, and share their experiences. Platforms such as Flipkart and Amazon generate enormous volumes of product reviews on a daily basis. These reviews serve as a primary source of information for prospective buyers, who rely on them to make informed purchasing decisions. However, with the exponential growth of online commerce, the sheer volume and diversity of user-generated content have made manual interpretation practically infeasible.

Sentiment Analysis, a key subfield of Natural Language Processing (NLP), refers to the computational process of identifying and categorizing the emotional tone expressed in a piece of text. By classifying text as positive, negative, or neutral, sentiment analysis tools provide automated insights into customer opinions, product reception, and overall market trends. With the increasing adoption of machine learning and deep learning techniques, modern sentiment analysis systems can handle the nuances of human language with remarkable accuracy.

Opinion Miner is a smart, interactive web system developed to address the challenge of analyzing online product reviews at scale. The system focuses on three complementary tasks: (1) detection and filtering of fake or deceptive reviews using a pre-trained Support Vector Machine (LinearSVC) with 91.55\% accuracy, (2) sentiment classification of genuine review text using VADER with custom negation handling, and (3) Aspect-Based Sentiment Analysis (ABSA) across nine product dimensions. It is designed to process both raw user-entered text and automatically scraped Flipkart product reviews in real time, making it an end-to-end solution for opinion mining in the e-commerce domain.

The system is further extended with a Chrome Browser Extension that integrates directly with Flipkart product pages, enabling users to obtain instant sentiment summaries without navigating away from their shopping session. The application is containerised using Docker and deployed on Render.com, making it publicly accessible without any local setup requirements.

\section{MOTIVATION}

The motivation behind developing Opinion Miner arises from several real-world challenges observed in the current e-commerce review ecosystem.

The existing approach to review analysis is largely manual or based on simplistic star-rating systems, both of which fail to capture the depth and complexity of customer sentiment. A critical issue in online review analysis is the prevalence of fake, incentivized, or spam reviews. Research has shown that a significant proportion of product reviews on major e-commerce platforms may be fabricated, either to artificially inflate a product's rating or to sabotage competitors. This manipulation of the review system erodes consumer trust and leads to poor purchasing decisions. Traditional sentiment tools do not account for review authenticity, thereby producing unreliable sentiment outcomes.

Furthermore, existing global sentiment tools are primarily trained on English text and fail to capture nuanced, domain-specific product vocabulary commonly found in Indian e-commerce reviews. The development of Opinion Miner is therefore driven by the need for a reliable, interactive, and comprehensive review analysis platform that combines real-time scraping, ML-powered fake review filtering, aspect-level sentiment analysis, and visual output within a single, user-friendly interface.

\section{PROBLEM DEFINITION}

The current state of product review analysis on e-commerce platforms is characterized by several key limitations that hinder consumers and businesses alike. The following problems have been identified and form the basis of this project:

\begin{itemize}
  \item Organizations and users face difficulty in understanding customer opinions due to the overwhelming volume of textual feedback. There is no scalable, automated mechanism available to the general public for interpreting hundreds of reviews for a single product.

  \item Manual analysis of product reviews is impractical and highly prone to bias. Human readers may selectively focus on certain reviews, leading to skewed conclusions about overall product quality.

  \item The pervasive presence of fake and deceptive reviews on platforms like Flipkart distorts sentiment analysis outcomes. Without a mechanism to detect and filter such reviews before analysis, the resulting sentiment scores are fundamentally unreliable.

  \item Existing tools provide only a single aggregate sentiment label without revealing which specific product attributes (such as build quality, delivery speed, or battery life) are driving positive or negative sentiment.

  \item Existing platforms do not support product-to-product sentiment comparison, which is a critical need for consumers choosing between multiple similar products.

  \item There is a lack of simple and interactive platforms that can analyze both user-input text and real-world product reviews effectively. Most available tools operate as either offline classifiers or inaccessible research prototypes, offering no practical user interface.
\end{itemize}

These problems collectively highlight the need for a unified, intelligent, and accessible review analysis system that goes beyond simple keyword matching or star-rating aggregation.

\section{OBJECTIVE OF THE PROJECT}

The primary objective of Opinion Miner is to develop a reliable, real-time, and explainable sentiment analysis system for online product reviews, with integrated fake review detection, aspect-level analysis, and visualisation capabilities. The specific objectives are listed below:

\begin{enumerate}
  \item To implement a VADER-based sentiment analysis engine capable of computing positive, negative, neutral, and compound sentiment scores for any input text, with custom negation handling to improve accuracy on review data.

  \item To develop a Selenium-powered Flipkart review scraper that can handle dynamic page loading using headless Chrome with webdriver-manager and extract real-time product reviews for sentiment analysis directly from a product URL.

  \item To build a Machine Learning-based Fake Review Detection module using a pre-trained LinearSVC Support Vector Machine with TF-IDF features (15,000 unigrams and bigrams) trained on the Cornell Deceptive Opinion Spam dataset of 40,432 reviews, achieving 91.55\% classification accuracy.

  \item To implement Aspect-Based Sentiment Analysis (ABSA) that decomposes product sentiment into nine distinct aspects---Quality, Delivery, Value, Performance, Battery, Camera, Display, Design, and Service---using sentence-level tokenisation and VADER scoring per aspect.

  \item To design and implement a Flask-based interactive web application with a clean, responsive user interface that allows users to input text or a Flipkart product URL and receive detailed sentiment analysis results.

  \item To visualize sentiment results using bar charts, score charts, and word clouds powered by Matplotlib and WordCloud libraries, making results easy to interpret for non-technical users.

  \item To develop a Chrome Browser Extension (Manifest V3) that integrates directly with e-commerce product pages, enabling users to obtain instant sentiment summaries without navigating away from their shopping session.

  \item To containerise the application using Docker and deploy it on Render.com for public cloud access, using Gunicorn as the WSGI server for production-grade request handling.

  \item To design a modular, maintainable codebase with clear separation of concerns across scraping, NLP, aggregation, and visualisation modules, enabling easy future extension.
\end{enumerate}

In summary, Opinion Miner aims to serve as a comprehensive, user-friendly, and technically robust platform that bridges the gap between raw e-commerce review data and actionable, trustworthy consumer insights.

% ══════════════════════════════════════════════════════════════════════════════
%  CHAPTER 2: ANALYSIS
% ══════════════════════════════════════════════════════════════════════════════
\chapter{Analysis}

\section{USER REQUIREMENT ANALYSIS}

The Opinion Miner system is designed as a multi-role, web-based sentiment analysis platform serving three primary user groups: General Users (consumers), System Administrators, and Developers or Researchers. To ensure that the platform meets functional expectations and performs reliably under varied conditions, both functional and non-functional requirements have been analyzed carefully before system design and implementation.

\subsection{Functional Requirements}

\subsubsection*{1.\quad General User Requirements}

\begin{itemize}
  \item Users should be able to enter free-form text through the manual review input box and receive instant sentiment analysis results including sentiment label (Positive, Negative, or Neutral) along with compound, positive, negative, and neutral scores.

  \item Users should be able to submit a valid Flipkart product review page URL, after which the system automatically scrapes available reviews, filters fake reviews using the SVM model, applies VADER sentiment analysis, performs aspect-based analysis, and returns an aggregated product sentiment summary.

  \item Users should be able to view a count-level breakdown of reviews categorised as Positive, Negative, Neutral, and Fake, displayed alongside the final product sentiment decision and average compound, positive, neutral, and negative scores.

  \item Users must be able to interact with the Chrome Browser Extension, which should detect when they are browsing a Flipkart product page and provide a sentiment analysis popup on demand with aspect breakdown.

  \item Users should receive visual feedback through bar charts, score charts, and word clouds for a comprehensive understanding of review patterns.
\end{itemize}

\subsubsection*{2.\quad Text Processing Requirements}

\begin{itemize}
  \item The system must preprocess all input text before sentiment analysis. Preprocessing steps must include lowercasing, punctuation removal, stopword removal, tokenization, and lemmatization.

  \item The preprocessing pipeline must handle special characters, URLs, HTML tags, and numeric content gracefully, stripping them from the text without corrupting meaningful tokens.

  \item Negation words (not, never, no, without, barely, hardly) must be retained in the token list as they are critical for the custom negation handling logic in the sentiment module.
\end{itemize}

\subsubsection*{3.\quad Sentiment Analysis Module Requirements}

\begin{itemize}
  \item The system must use the VADER model as the primary sentiment analysis engine due to its suitability for short, informal, and social-media-style text common in product reviews.

  \item The sentiment module must apply custom negation handling logic to adjust compound scores when negation words are detected in the review.

  \item The system must compute aggregated sentiment scores (average compound, positive, neutral, negative) across all genuine reviews for any given Flipkart product.

  \item The pipeline must clearly separate fake review filtering from sentiment computation: only reviews classified as genuine by the SVM must be passed to the sentiment analysis engine.
\end{itemize}

\subsubsection*{4.\quad Fake Review Detection Requirements}

\begin{itemize}
  \item The fake review detection module must use a pre-trained LinearSVC model with TF-IDF vectorisation to classify each scraped review as Genuine (OR) or Fake (CG).

  \item Reviews classified as Fake must be discarded from the sentiment analysis pipeline and reported separately in the results summary with their predicted label and detection method.

  \item The system must display the count of fake reviews detected for any given product analysis.
\end{itemize}

\subsubsection*{5.\quad Web Scraping Module Requirements}

\begin{itemize}
  \item The system must use a Selenium-based scraper configured with ChromeDriver (auto-managed via webdriver-manager) to navigate Flipkart product review pages and extract review text, star ratings, reviewer names, and review dates.

  \item The scraper must skip all Snap-packaged Chromium binaries and prefer system-installed Google Chrome to avoid headless mode incompatibility.

  \item The scraper must implement fallback CSS selectors to account for periodic changes in Flipkart's React Native Web (RNW) page structure, ensuring robustness against class name alterations.

  \item The scraper must support dynamic page navigation (up to 50 pages per product) to load additional reviews beyond the initially visible set.
\end{itemize}

\subsubsection*{6.\quad System-Level (Cross-Functional) Requirements}

\begin{itemize}
  \item The system must provide a Flask-based REST API backend that accepts inputs from both the web interface and the Chrome Extension, processes them, and returns JSON-formatted results.

  \item All results must be returned within an acceptable response time, targeting under ten seconds for text analysis and under ninety seconds for full Flipkart product URL analysis (50 pages).

  \item The frontend interface must be responsive and functional across major modern browsers including Chrome, Firefox, Edge, and Safari.
\end{itemize}

\subsection{Non-Functional Requirements}

\begin{itemize}
  \item \textbf{Performance:} The system must process a single text input and return sentiment results in real time, within two to three seconds. For product URL analysis involving scraping of up to 50 review pages, total response time should remain under ninety seconds under normal network conditions.

  \item \textbf{Usability:} The interface must be intuitive, with clearly labeled input fields, result panels, aspect cards, and visual charts. Non-technical users should be able to perform full sentiment analysis without requiring any prior knowledge of NLP.

  \item \textbf{Reliability:} The system must handle scraping failures gracefully, displaying a meaningful error message if a product URL is invalid, the page is blocked, or no reviews are found, rather than crashing or hanging indefinitely.

  \item \textbf{Scalability:} The system architecture should be modular, allowing individual components such as the scraping module, sentiment engine, and fake detection module to be independently updated or replaced without affecting the overall pipeline.

  \item \textbf{Maintainability:} All code must follow PEP 8 style guidelines for Python and be organised into clearly named modules with inline documentation, enabling easy debugging and future feature additions.

  \item \textbf{Compatibility:} The application must run on Windows, Linux, and macOS operating systems. The Docker deployment ensures consistent behaviour across environments. The Chrome Extension must be compatible with Google Chrome version 90 and above (Manifest V3).

  \item \textbf{Security:} The system must not store any user-entered review text or product URLs beyond the duration of an active session. No personal data collection or tracking must be implemented.
\end{itemize}

\subsection{Hardware and Software Requirements}

\subsubsection*{Hardware Requirements}

\begin{itemize}
  \item Processor: Intel Core i3 (4th Generation) or higher, or equivalent AMD processor
  \item RAM: Minimum 8 GB (16 GB recommended for running Selenium alongside Flask server)
  \item Storage: Minimum 512 GB HDD or 128 GB SSD
  \item Internet Connection: Required for Flipkart URL scraping and real-time review fetching
  \item Display: Minimum resolution of 1280 $\times$ 720 for full interface visibility
\end{itemize}

\subsubsection*{Software Requirements}

\begin{itemize}
  \item Operating System: Windows 10 / Ubuntu 20.04 LTS / macOS Big Sur or later
  \item Programming Language: Python 3.8 or above
  \item Web Framework: Flask 2.x
  \item NLP Libraries: NLTK 3.x, vaderSentiment
  \item Machine Learning: scikit-learn (LinearSVC, TfidfVectorizer)
  \item Web Scraping: Selenium 4.x with webdriver-manager (auto ChromeDriver matching)
  \item Data Visualization: Matplotlib 3.x, WordCloud
  \item Frontend: HTML5, CSS3, JavaScript (ES6+)
  \item Browser Extension: Chrome Extension Manifest V3
  \item Deployment: Docker, Gunicorn, Render.com
  \item Development Environment: Visual Studio Code, Jupyter Notebook
  \item Version Control: Git and GitHub
\end{itemize}

% ══════════════════════════════════════════════════════════════════════════════
%  CHAPTER 3: SYSTEM ANALYSIS
% ══════════════════════════════════════════════════════════════════════════════
\chapter{System Analysis}

\section{EXISTING SYSTEM}

The current landscape of online product review analysis and consumer sentiment understanding relies predominantly on manual or primitive automated methods, both of which present significant limitations in the context of modern e-commerce scale and diversity.

The most widely practiced existing approach involves consumers manually reading through product reviews listed on e-commerce platforms. This is time-consuming and cognitively demanding, particularly for products that have accumulated hundreds or thousands of reviews. Consumers who wish to gauge the overall reception of a product must invest considerable effort in sifting through reviews without any structured summary of the collective sentiment. This approach is also highly subjective, as different readers may interpret the same review differently.

Platforms like Flipkart and Amazon provide star-rating-based decision aids, displaying an aggregate rating score alongside a count of ratings at each star level. While this provides a quick visual overview, it fails to capture the rich semantic content of written reviews. A product with a 4.0-star average may receive that rating from both highly satisfied customers and moderately disappointed ones, making the rating alone an unreliable indicator of product quality for specific use cases.

Some basic sentiment tools and browser extensions exist that perform rudimentary keyword-based analysis of review text. These tools classify sentiment based on the presence of predefined positive or negative words without understanding sentence context, negation, or sentiment intensity. As a result, they exhibit poor accuracy for complex sentences and are unable to differentiate between sarcastic, exaggerated, or domain-specific product text.

The major challenges and drawbacks identified in existing systems include the following:

\begin{itemize}
  \item \textbf{Time Consuming:} Manual review reading across hundreds of product listings is impractical, especially for frequent online shoppers or business analysts who need to compare multiple competing products.

  \item \textbf{No Explanation for Sentiment:} Existing basic tools output only a sentiment label without providing the underlying scores, contributing factors, or justification for the classification, making results opaque and difficult to trust.

  \item \textbf{No Fake Review Filtering:} None of the commonly used consumer-facing tools incorporate a mechanism to detect or exclude fake, incentivized, or spammy reviews from the analysis. Sentiment scores derived from manipulated reviews are inherently misleading.

  \item \textbf{No Aspect-Level Insights:} Existing platforms provide only a single overall sentiment score without revealing which specific product dimensions (Quality, Delivery, Battery, Camera, etc.) customers are praising or criticising.

  \item \textbf{Poor Accuracy for Complex Sentences:} Keyword-based approaches cannot handle negation (e.g., `not bad' being misclassified as negative), sarcasm, comparative expressions, or domain-specific jargon, leading to frequent misclassification.

  \item \textbf{No Product-to-Product Comparison:} Existing tools do not support side-by-side sentiment comparison of multiple products, which is a critical requirement for purchase-decision support in competitive product categories.
\end{itemize}

Thus, the existing landscape is fragmented, inaccurate, and unable to meet the demands of informed, real-time, and trustworthy e-commerce review analysis.

\section{PROPOSED SYSTEM}

The proposed system, Opinion Miner, introduces a comprehensive, intelligent, and user-friendly solution for analysing online product reviews in real time. It addresses every identified limitation of the existing system through a modular, technology-driven architecture that integrates automated scraping, text preprocessing, ML-powered fake review detection, VADER sentiment analysis, and aspect-based sentiment analysis within a single interactive platform.

The key features of the proposed system are described below:

\begin{itemize}
  \item \textbf{Real-Time Flipkart Review Scraping:} The system uses a Selenium-based web scraper to extract product reviews directly from any valid Flipkart product URL. The scraper handles dynamic page loading through multi-page navigation (up to 50 pages) and uses webdriver-manager to automatically pair the correct ChromeDriver version with the installed Google Chrome binary, bypassing Snap-path incompatibilities.

  \item \textbf{Automated Text Preprocessing Pipeline:} All review text is passed through a preprocessing pipeline that performs lowercasing, punctuation removal, stopword filtering (retaining negation words), tokenisation, and lemmatisation using NLTK's WordNetLemmatizer before sentiment analysis.

  \item \textbf{ML-Based Fake Review Detection:} The fake review detection module employs a pre-trained LinearSVC Support Vector Machine with a TF-IDF vectoriser (15,000 features, unigrams and bigrams). Trained on 40,432 labelled reviews (20,216 genuine OR and 20,216 fake CG) from the Cornell Deceptive Opinion Spam dataset, the model achieves 91.55\% accuracy, 91.63\% precision, 91.47\% recall, and 91.55\% F1-score. Only reviews classified as genuine are passed to the sentiment pipeline.

  \item \textbf{VADER Sentiment Analysis with Custom Negation Handling:} The system employs the VADER sentiment model, well-suited for short, informal, and social media-style text, to compute positive, negative, neutral, and compound sentiment scores. A custom negation handling function adjusts compound scores when negation words (not, never, no, without, barely, hardly) are detected in the review, improving overall classification accuracy. Thresholds: compound $> 0.2$ = Positive; compound $< -0.2$ = Negative; otherwise Neutral.

  \item \textbf{Aspect-Based Sentiment Analysis (ABSA):} Reviews are split into sentences using NLTK's \texttt{sent\_tokenize}. Each sentence is matched against keyword lists for nine product aspects: Quality, Delivery, Value, Performance, Battery, Camera, Display, Design, and Service. VADER is applied per matched sentence, and per-aspect average compound scores, verdict labels, and sample sentences are returned.

  \item \textbf{Aggregated Product Sentiment:} For Flipkart product analysis, the system aggregates the sentiment scores of all verified genuine reviews to produce a final product-level verdict of Positive, Negative, or Neutral, along with count-level breakdowns and average scores (compound, positive, neutral, negative).

  \item \textbf{Interactive Web Interface:} The Flask-based web application provides a clean, dark-themed responsive interface with separate input panels for text analysis and URL-based Flipkart analysis. Results are displayed with sentiment labels, score summaries, aspect cards with colour-coded bars, and embedded chart visualisations.

  \item \textbf{Chrome Browser Extension:} A browser extension (Manifest V3) integrates with Flipkart product pages, extracts the current product URL, sends it to the Flask backend for real-time analysis, and displays results in a popup panel with aspect breakdown.

  \item \textbf{Docker Deployment:} The application is fully containerised using Docker with Chromium pre-installed, and deployed on Render.com using Gunicorn as the production WSGI server for reliable, scalable public access.
\end{itemize}

The advantages of the proposed system over existing approaches are summarized below:

\begin{itemize}
  \item Automates the complete review analysis workflow from data collection to insight generation.
  \item Eliminates the influence of fake reviews on sentiment outcomes with 91.55\% ML accuracy.
  \item Provides granular aspect-level insights across 9 product dimensions.
  \item Provides explainable results with score breakdowns and visual charts, making outputs interpretable for non-technical users.
  \item Operates as a live, browser-integrated tool via the Chrome Extension, enabling seamless use during the shopping experience.
  \item Production-ready with Docker containerisation and cloud deployment on Render.com.
\end{itemize}

\section{SOFTWARE USED}

The following technologies and tools were used in the development of Opinion Miner:

\subsubsection*{Core Programming Technologies}

\begin{itemize}
  \item \textbf{Python (Version 3.11):} The primary programming language used for backend development, NLP processing, machine learning, web scraping, and server logic. Python's extensive library ecosystem for data science and NLP made it the natural choice for this project.

  \item \textbf{Jupyter Notebook:} Used during the experimentation and prototyping phase for developing, testing, and iterating on the text preprocessing pipeline, VADER sentiment implementation, and SVM fake review detection logic before integration into the Flask application.
\end{itemize}

\subsubsection*{Web Framework and Backend}

\begin{itemize}
  \item \textbf{Flask (Version 2.x):} A lightweight Python-based web framework used to build the backend REST API server. Flask's simplicity and flexibility made it suitable for routing user requests between the frontend interface and the NLP pipeline modules.

  \item \textbf{Gunicorn:} A production-grade Python WSGI HTTP server used to serve the Flask application in the Docker container with a 600-second timeout for long scraping operations.
\end{itemize}

\subsubsection*{NLP and Machine Learning Libraries}

\begin{itemize}
  \item \textbf{NLTK (Natural Language Toolkit):} Used for standard text preprocessing operations including tokenisation, stopword removal, and lemmatisation through the WordNetLemmatizer utility.

  \item \textbf{vaderSentiment:} The primary sentiment analysis library, providing the SentimentIntensityAnalyzer class. VADER is specifically calibrated for social media and short-form text, making it well-suited for e-commerce reviews.

  \item \textbf{scikit-learn:} Used for the fake review detection model. Specifically, \texttt{LinearSVC} for the SVM classifier and \texttt{TfidfVectorizer} for feature extraction with 15,000 features and bigram support.
\end{itemize}

\subsubsection*{Web Scraping}

\begin{itemize}
  \item \textbf{Selenium (Version 4.x) with webdriver-manager:} Used to automate browser interaction with Flipkart product pages. webdriver-manager automatically downloads and manages the correct ChromeDriver version matching the installed Google Chrome binary, resolving version incompatibility issues.

  \item \textbf{BeautifulSoup4:} Used as a supplementary HTML parsing library to extract structured content from the page source once Selenium has fully rendered the target review page.
\end{itemize}

\subsubsection*{Data Visualisation}

\begin{itemize}
  \item \textbf{Matplotlib (Version 3.x):} Used to generate static bar charts representing sentiment distribution and score comparisons across product reviews. Charts are encoded as Base64 PNG strings for inline display in the frontend.

  \item \textbf{WordCloud:} Used to produce positive and negative word cloud visualisations from the most frequently occurring terms in genuine positive and negative reviews respectively.
\end{itemize}

\subsubsection*{Deployment and DevOps}

\begin{itemize}
  \item \textbf{Docker:} Used to containerise the entire application including Python 3.11-slim base, Chromium browser, ChromeDriver, and all Python dependencies. Environment variables \texttt{CHROME\_BIN} and \texttt{CHROMEDRIVER\_BIN} configure the scraper to use the correct binaries inside the container.

  \item \textbf{Render.com:} Cloud hosting platform used for public deployment of the Docker container. The \texttt{render.yaml} configuration specifies the web service, runtime, health check path, and environment variables.
\end{itemize}

\subsubsection*{Frontend Technologies}

\begin{itemize}
  \item \textbf{HTML5, CSS3, JavaScript (ES6+):} Used to build the dark-themed web interface including input forms, result panels, aspect cards with colour-coded bars, and dynamic chart rendering using asynchronous AJAX calls to the Flask backend.
\end{itemize}

\subsubsection*{Development Tools}

\begin{itemize}
  \item \textbf{Visual Studio Code:} Primary integrated development environment for writing, debugging, and managing the application codebase.

  \item \textbf{Git and GitHub:} For version control, team collaboration, and codebase management throughout the development lifecycle. Repository: \textit{ThagillaRakesh/Flipkart-Sentiment-Analysis}.
\end{itemize}

% ══════════════════════════════════════════════════════════════════════════════
%  CHAPTER 4: SYSTEM DESIGN
% ══════════════════════════════════════════════════════════════════════════════
\chapter{System Design}

System design constitutes the blueprint of the proposed system. It explains the overall structure, the interaction between system components, and the movement of data across modules. In Opinion Miner, system design encompasses use case modeling, class diagrams, sequence diagrams, and data flow diagrams. These collectively help in understanding the workflow and the dependencies between different user roles and functional modules.

\section{UML DIAGRAMS}

Unified Modeling Language (UML) diagrams are used to describe, visualise, and document the structure and behaviour of the system. They represent how different entities and components interact within the architecture of Opinion Miner.

\subsubsection*{1.\quad Use Case Diagram}

The Use Case Diagram for Opinion Miner illustrates the interaction between the three primary actor types and the system. The \textbf{General User} can: perform text sentiment analysis, initiate Flipkart product review analysis through URL submission, interact with the Chrome Extension, view visualisation results (bar charts, word clouds), and view aspect-based sentiment breakdown. The \textbf{System Administrator} can: monitor system activity, manage deployment configuration, and update model artifacts. The \textbf{Browser Extension User}, a specialised user role, interacts with the system exclusively through the Chrome Extension popup panel on Flipkart product pages.

\begin{figure}[H]
\centering
\fbox{\parbox{12cm}{\centering\vspace{3cm}[ Place Figure 4.1: Use Case Diagram here ]\vspace{3cm}}}
\caption{Use Case Diagram}
\label{fig:usecase}
\end{figure}

\subsubsection*{2.\quad Class Diagram}

The Class Diagram represents the primary entities within the Opinion Miner system and their relationships. The main classes include: \textbf{ReviewInput} (encapsulating raw text and URL inputs), \textbf{Preprocessor} (with methods for cleaning, tokenising, and normalising text including lemmatisation), \textbf{FakeReviewDetector} (wrapping the LinearSVC model with TF-IDF vectoriser, with a \texttt{classify()} method), \textbf{SentimentAnalyser} (wrapping the VADER model with custom negation logic), \textbf{AspectAnalyser} (performing sentence-level ABSA across 9 aspects), \textbf{Scraper} (managing Selenium browser sessions for Flipkart review extraction), \textbf{ResultAggregator} (computing product-level sentiment summaries and average scores), and \textbf{VisualisationModule} (generating charts and word clouds from sentiment data).

\begin{figure}[H]
\centering
\fbox{\parbox{12cm}{\centering\vspace{3cm}[ Place Figure 4.2: Class Diagram here ]\vspace{3cm}}}
\caption{Class Diagram}
\label{fig:class}
\end{figure}

\subsubsection*{3.\quad Sequence Diagram}

The Sequence Diagram demonstrates the order of interactions between system components during a typical Flipkart product analysis operation. The sequence begins with the user submitting a product URL through the web interface. The Flask server receives the request and triggers the Selenium Scraper, which navigates to the product page, waits for content to load, scrolls through up to 50 pages, and extracts review text along with star ratings. Each extracted review is passed to the FakeReviewDetector, which applies the LinearSVC model. Reviews classified as genuine are forwarded to the Preprocessor and subsequently to the SentimentAnalyser. The AspectAnalyser performs sentence-level analysis on each genuine review. The ResultAggregator compiles individual scores into a product-level summary, which is formatted as a JSON response and returned to the frontend for display.

\begin{figure}[H]
\centering
\fbox{\parbox{12cm}{\centering\vspace{3cm}[ Place Figure 4.3: Sequence Diagram here ]\vspace{3cm}}}
\caption{Sequence Diagram}
\label{fig:sequence}
\end{figure}

\section{DATA FLOW}

The Data Flow Diagram (DFD) illustrates how information enters, is processed, and exits the Opinion Miner system at different levels of abstraction.

\subsubsection*{1.\quad DFD Context Level (Level 0)}

At the context level, Opinion Miner is represented as a single process receiving inputs from two external entities: the \textbf{User} (via the web interface or Chrome Extension) and the \textbf{Flipkart Website} (as a data source for scraped reviews). The outputs produced by the system include Sentiment Results (labels, scores, and visualisations), Fake Review Reports (count and detection method), Aspect Sentiment Cards (per-aspect verdict and avg\_compound), and Product Sentiment Summaries. The system interacts bidirectionally with the User, receiving either raw text or product URLs as input and returning analysis results.

\begin{figure}[H]
\centering
\fbox{\parbox{12cm}{\centering\vspace{3cm}[ Place Figure 4.4: Context Level DFD here ]\vspace{3cm}}}
\caption{Context Level DFD}
\label{fig:dfd}
\end{figure}

\subsubsection*{2.\quad Explanation of Data Flow}

The detailed data flow within the system proceeds as follows. When a user submits a Flipkart product URL, the input is received by the Flask backend and routed to the Scraper module. The Scraper launches a headless Chrome browser session via Selenium, navigates to the product review page, extracts review text and ratings across up to 50 pages, and passes them to the FakeReviewDetector. Reviews passing the ML authenticity check (classified as genuine by LinearSVC) are forwarded to the Text Preprocessor, which cleans and normalises the text. The preprocessed text is then fed into the VADER SentimentAnalyser, which produces individual review sentiment scores. These scores are used by the AspectAnalyser to perform sentence-level ABSA. The ResultAggregator compiles counts, averages, and overall verdicts. The VisualisationModule generates charts encoded as Base64 PNG. The final JSON response, comprising the product verdict, score distribution, aspect cards, and visualisation data, is returned to the frontend and rendered in the browser.

For text input analysis, the flow is simpler: the user-entered text bypasses the scraping stage and enters the pipeline directly at the SentimentAnalyser, following preprocessing. The Chrome Extension sends the current Flipkart product URL directly to the same Flask API endpoint, enabling it to leverage the full analysis pipeline.

\subsubsection*{3.\quad Data Storage}

Opinion Miner is designed as a stateless, session-based system. No persistent database is used for storing user inputs or analysis results. Temporary data, including scraped review text and computed sentiment scores, exists in memory only for the duration of an active request-response cycle. The pre-trained SVM model (\texttt{svm\_model.pkl}) and TF-IDF vectoriser (\texttt{tfidf\_svm.pkl}) are loaded once at application startup and cached in memory for efficient inference.

% ══════════════════════════════════════════════════════════════════════════════
%  CHAPTER 5: IMPLEMENTATION
% ══════════════════════════════════════════════════════════════════════════════
\chapter{Implementation}

Implementation refers to the actual construction and integration of the software system as per the design and specifications. Opinion Miner is implemented using Python as the core language, with Flask for the web backend, Selenium for real-time review scraping, NLTK and vaderSentiment for NLP processing, scikit-learn for ML-based fake detection, and HTML/CSS/JavaScript for the frontend interface. This chapter details the project structure, implementation methodology, module descriptions, and representative screenshots of the running application.

\section{PROJECT STRUCTURE}

The entire system is organised into a modular directory structure separating frontend assets, backend logic, NLP modules, scraping utilities, and configuration files. The structure is as follows:

\begin{singlespace}
\begin{lstlisting}
opinion_miner/
|-- app.py                    (Flask application entry point)
|-- requirements.txt          (Python dependency list)
|-- Dockerfile                (Docker container definition)
|-- render.yaml               (Render.com deployment config)
|-- train_svm.py              (SVM model training script)
|-- scraper/
|   |-- flipkart_scraper.py   (Selenium-based review scraper)
|-- nlp/
|   |-- sentiment.py          (VADER engine with negation handler)
|   |-- fake_detector.py      (LinearSVC SVM fake review detector)
|   |-- aspect_sentiment.py   (Aspect-based sentiment analysis)
|   `-- models/
|       |-- svm_model.pkl     (Trained LinearSVC model)
|       `-- tfidf_svm.pkl     (Fitted TF-IDF vectoriser)
|-- utils/
|   |-- aggregator.py         (Review aggregation and scoring)
|   `-- visualizer.py         (Chart and word cloud generation)
|-- templates/
|   `-- index.html            (Main web interface)
|-- static/
|   |-- css/  style.css       (Frontend styling)
|   `-- js/   main.js         (Frontend interaction logic)
`-- chrome_extension/
    |-- manifest.json         (Chrome Extension manifest V3)
    |-- popup.html            (Extension popup interface)
    |-- popup.js              (Extension logic)
    `-- popup.css             (Extension styling)
\end{lstlisting}
\end{singlespace}

This modular folder structure ensures clear separation of concerns, enabling easy debugging, testing, and future extension of individual components without disrupting the overall system.

\section{DESCRIPTION OF IMPLEMENTATION}

\subsubsection*{Frontend Implementation}

The frontend of Opinion Miner is built using HTML5, CSS3, and vanilla JavaScript. The main page presents two distinct input panels: one for manual text entry and one for Flipkart product URL submission. Both panels include submit buttons that trigger asynchronous AJAX requests to the Flask backend, preventing full page reloads and providing a smooth user experience. Results are dynamically rendered in a dedicated results section below the input panels, displaying the sentiment label, numerical scores (compound, positive, neutral, negative averages), a review breakdown table, aspect cards with colour-coded bars, and embedded charts.

\subsubsection*{Backend Implementation}

The backend is implemented using the Flask micro-framework in Python. The application defines two primary API routes: \texttt{/analyze} for text input analysis and \texttt{/analyze-url} for Flipkart product URL analysis. Each route receives the input via an HTTP POST request, validates it, and dispatches it through the appropriate processing pipeline. The Flask server handles cross-origin requests from the Chrome Extension using appropriate CORS headers. All responses are returned as JSON objects containing the sentiment label, score dictionary, review breakdown statistics, aspect analysis data, platform badge, and Base64-encoded chart images for inline display in the frontend.

\subsubsection*{NLP Pipeline Implementation}

The NLP pipeline is implemented as a sequence of modular Python functions. The Preprocessor module accepts raw review text and performs: conversion to lowercase, removal of URLs and HTML tags using regex, stripping of punctuation characters, tokenisation using NLTK's \texttt{word\_tokenize} function, removal of English stopwords (retaining negation words), and lemmatisation using \texttt{WordNetLemmatizer}. The cleaned token list is reassembled into a normalised string for input to the VADER analyser.

The VADER \texttt{SentimentIntensityAnalyzer} is initialised once at application startup and reused across requests for efficiency. A custom negation check is applied: if any token in the review matches entries from a predefined negation word list (not, never, no, without, barely, hardly, scarcely, doesn't, isn't, wasn't, shouldn't, wouldn't, couldn't, won't, can't), the compound score is scaled by a dampening factor of 0.85 to reflect the reduced sentiment intensity introduced by negation.

\subsubsection*{Scraper Implementation}

The Selenium scraper initialises a Chrome WebDriver in headless mode (\texttt{--headless=new}) to avoid opening a visible browser window. The scraper uses webdriver-manager's \texttt{ChromeDriverManager().install()} to automatically download and cache the correct ChromeDriver version matching the installed Google Chrome. All Snap-packaged Chromium paths (containing \texttt{/snap/}) are explicitly skipped to avoid sandbox incompatibilities. The scraper navigates Flipkart review pages iteratively (up to 50 pages), using \texttt{WebDriverWait} and \texttt{ExpectedConditions} for reliable page load handling. Review elements are located using the primary CSS selector \texttt{span.css-1jxf684} for Flipkart's React Native Web (RNW) layout. Extracted reviews are passed through junk filtering (address detection, minimum length, exact duplicate removal) before being returned to the main pipeline.

\section{MODULE DESCRIPTION}

The system is divided into six primary functional modules, each responsible for a distinct stage of the opinion mining pipeline.

\subsection{User Input Module}

\textbf{Purpose:} Collects and validates input data from the user through the web interface or the Chrome Extension.

\textbf{Features:}
\begin{itemize}
  \item Accepts raw review text entered directly by the user in the text analysis panel.
  \item Accepts a Flipkart product review page URL for automated scraping and analysis.
  \item Validates URL format before passing to the scraping module to prevent malformed requests.
  \item Provides clear error messages for empty inputs, invalid URLs, or server-side processing failures.
\end{itemize}

\textbf{Implementation Highlights:} The Flask route functions perform basic input validation including null checks and URL scheme verification. Input text is sanitized to remove leading and trailing whitespace before being forwarded to the preprocessing module.

\subsection{Review Scraping Module}

\textbf{Purpose:} Retrieves product reviews from live Flipkart product pages in real time.

\textbf{Features:}
\begin{itemize}
  \item Navigates to Flipkart product review pages using a Selenium-controlled headless Chrome browser managed by webdriver-manager.
  \item Implements iterative page navigation across up to 50 review pages to collect a statistically meaningful sample.
  \item Uses primary CSS selector \texttt{span.css-1jxf684} targeting Flipkart's React Native Web (RNW) text nodes, with junk filtering for address-pattern reviews, reviews under 4 words, and exact duplicates.
  \item Extracts review text and star ratings for downstream processing.
  \item Skips Snap-packaged Chromium paths to avoid headless mode sandbox incompatibilities.
\end{itemize}

\textbf{Implementation Highlights:} \texttt{WebDriverWait} with a ten-second timeout replaces hardcoded sleep calls for more reliable page load handling. The scraper module gracefully handles \texttt{StaleElementReferenceException} and \texttt{NoSuchElementException} by logging warnings and continuing with available data.

\subsection{Fake Review Detection Module}

\textbf{Purpose:} Identifies and filters deceptive, incentivized, or spam reviews before they can influence sentiment analysis results.

\textbf{Features:}
\begin{itemize}
  \item Uses a pre-trained \textbf{LinearSVC} Support Vector Machine model to classify each review as Genuine (OR) or Fake (CG).
  \item Feature extraction via TF-IDF vectoriser with 15,000 features, unigram and bigram range, minimum document frequency of 2, and sublinear TF scaling.
  \item Preprocessing pipeline: lowercasing, URL removal, non-alpha stripping, word tokenisation, stopword removal (retaining negation words), and WordNetLemmatizer lemmatisation.
  \item Trained on 40,432 balanced reviews (20,216 genuine OR + 20,216 fake CG) from the Cornell Deceptive Opinion Spam dataset.
  \item \textbf{Performance:} Accuracy 91.55\%, Precision 91.63\%, Recall 91.47\%, F1-Score 91.55\% (80/20 stratified split).
  \item Reviews classified as Fake are excluded from the sentiment pipeline and reported separately in the results summary with their label and detection method.
\end{itemize}

\subsection{Sentiment Analysis Module}

\textbf{Purpose:} Determines the emotional tone of each genuine review using the VADER sentiment model and custom negation logic.

\textbf{Features:}
\begin{itemize}
  \item Computes positive, negative, neutral, and compound sentiment scores for each review using the VADER \texttt{SentimentIntensityAnalyzer}.
  \item Applies custom negation handling: if negation words are detected, the compound score is reduced by a dampening factor of 0.85.
  \item Classifies each review as Positive (compound $> 0.2$), Negative (compound $< -0.2$), or Neutral (otherwise).
  \item Aggregates individual review scores to compute mean positive, negative, neutral, and compound scores for the full product review set.
  \item Produces a product-level sentiment verdict based on the distribution of individual review classifications.
\end{itemize}

\subsection{Aspect-Based Sentiment Analysis Module}

\textbf{Purpose:} Provides granular sentiment breakdown across nine product dimensions.

\textbf{Features:}
\begin{itemize}
  \item Reviews are split into sentences using NLTK's \texttt{sent\_tokenize}.
  \item Each sentence is matched against keyword lists for nine product aspects: \textbf{Quality, Delivery, Value, Performance, Battery, Camera, Display, Design,} and \textbf{Service}.
  \item VADER is applied per matched sentence to compute aspect-level sentiment scores.
  \item Returns per-aspect: positive/neutral/negative counts, average compound score, verdict label, and up to three sample sentences.
  \item Results are displayed as colour-coded aspect cards in the frontend with a horizontal bar showing the compound score.
\end{itemize}

\subsection{Visualisation Module}

\textbf{Purpose:} Presents the analysis results visually for intuitive interpretation.

\textbf{Features:}
\begin{itemize}
  \item Generates a \textbf{Sentiment Distribution Bar Chart} showing counts of Positive, Neutral, Negative, and Fake reviews using Matplotlib.
  \item Generates a \textbf{Sentiment Score Chart} showing average positive, neutral, negative, and compound scores.
  \item Produces a \textbf{Positive Word Cloud} from the most frequently occurring terms in positive reviews.
  \item Produces a \textbf{Negative Word Cloud} from the most frequently occurring terms in negative reviews.
  \item All charts are encoded as Base64 PNG strings and transmitted in the JSON API response for inline rendering in the browser.
\end{itemize}

\section{SCREENSHOTS}

The following figures represent the user interface and output screens of the Opinion Miner platform.

\begin{figure}[H]
\centering
\fbox{\parbox{12cm}{\centering\vspace{3cm}[ Place Figure 5.1: System Architecture Overview here ]\vspace{3cm}}}
\caption{System Architecture Overview}
\label{fig:arch}
\end{figure}

\begin{figure}[H]
\centering
\fbox{\parbox{12cm}{\centering\vspace{3cm}[ Place Figure 5.2: Text Sentiment Analysis Interface here ]\vspace{3cm}}}
\caption{Text Sentiment Analysis Interface}
\label{fig:textui}
\end{figure}

\begin{figure}[H]
\centering
\fbox{\parbox{12cm}{\centering\vspace{3cm}[ Place Figure 5.3: Flipkart Review Sentiment Interface here ]\vspace{3cm}}}
\caption{Flipkart Review Analysis Interface}
\label{fig:urlui}
\end{figure}

\begin{figure}[H]
\centering
\fbox{\parbox{12cm}{\centering\vspace{3cm}[ Place Figure 5.4: Sentiment Result and Visualisation Output here ]\vspace{3cm}}}
\caption{Sentiment Result and Visualisation Output}
\label{fig:results}
\end{figure}

\begin{figure}[H]
\centering
\fbox{\parbox{12cm}{\centering\vspace{3cm}[ Place Figure 5.5: Aspect-Based Sentiment Cards here ]\vspace{3cm}}}
\caption{Aspect-Based Sentiment Analysis Cards}
\label{fig:aspects}
\end{figure}

\begin{figure}[H]
\centering
\fbox{\parbox{12cm}{\centering\vspace{3cm}[ Place Figure 5.6: Chrome Extension Popup here ]\vspace{3cm}}}
\caption{Chrome Extension Popup Interface}
\label{fig:extension}
\end{figure}

% ══════════════════════════════════════════════════════════════════════════════
%  CHAPTER 6: TESTING
% ══════════════════════════════════════════════════════════════════════════════
\chapter{Testing}

Software testing is a critical phase in the development lifecycle that ensures the reliability, accuracy, and stability of the application. For Opinion Miner, testing was conducted at multiple levels including system testing, unit testing, security testing, and integration testing to verify that all modules function correctly and interact seamlessly as a complete system.

\section{SYSTEM TESTING}

System testing validates the complete and integrated Opinion Miner application to ensure that all functionalities perform as intended. It focuses on verifying both functional and non-functional requirements under realistic usage conditions.

\subsubsection*{Objectives of System Testing:}

\begin{itemize}
  \item Ensure the complete Opinion Miner platform meets all specified user requirements including text analysis, URL-based Flipkart analysis, ML fake review detection, aspect-based sentiment analysis, and visualisation.
  \item Verify that all modules---Scraping, Preprocessing, Sentiment Analysis, Fake Detection, Aspect Analysis, and Visualisation---interact without conflicts.
  \item Validate end-to-end flows such as URL submission, scraping, fake filtering, sentiment computation, aspect scoring, and result display.
  \item Confirm proper data flow between the frontend, Flask backend, and NLP pipeline modules.
\end{itemize}

\subsubsection*{System Testing Scenarios and Results:}

\begin{table}[H]
\centering
\begin{tabular}{|p{3.5cm}|p{4cm}|p{4cm}|p{1.8cm}|}
\hline
\textbf{Test Scenario} & \textbf{Expected Result} & \textbf{Actual Result} & \textbf{Status} \\
\hline
Text Input Analysis & Sentiment label and scores returned & Label and scores displayed correctly & Pass \\
\hline
Flipkart URL Scraping & Reviews extracted from URL & Reviews extracted and listed & Pass \\
\hline
SVM Fake Review Detection & Fake reviews classified and excluded & Fake reviews filtered correctly & Pass \\
\hline
Aspect-Based Analysis & 9 aspect cards with verdicts shown & Aspect cards rendered correctly & Pass \\
\hline
Negation Handling & Compound score adjusted for `not bad' & Score reduced by 0.85x factor & Pass \\
\hline
Average Scores Display & Avg compound/pos/neu/neg shown & Scores displayed in verdict box & Pass \\
\hline
Empty Input Submission & Error message displayed & Error message rendered & Pass \\
\hline
Invalid URL Submission & Error message for bad URL & Validation error shown & Pass \\
\hline
Sentiment Visualisation & Bar chart and word clouds generated & Charts rendered in results panel & Pass \\
\hline
Chrome Extension Analysis & Popup displays sentiment results & Extension popup working & Pass \\
\hline
\end{tabular}
\caption{System Testing Scenarios and Results}
\end{table}

System testing confirmed that Opinion Miner integrates all modules with acceptable latency, providing accurate responses and reliable behaviour under normal operating conditions.

\section{UNIT TESTING}

Unit testing involves verifying individual modules or components of the system in isolation to ensure they perform their specific functions correctly and independently of other modules.

\subsubsection*{Sample Unit Test Cases:}

\begin{itemize}
  \item \textbf{Preprocessor Module:} Verified that the preprocessing function correctly lowercases input, removes URLs, strips punctuation, filters stopwords, retains negation words (not, never, no), and returns a non-empty normalised string for a valid review input. Verified that an empty string input returns an empty string without raising exceptions.

  \item \textbf{VADER Sentiment Module:} Verified that \texttt{analyze()} returns a dictionary with keys `pos', `neg', `neu', and `compound' for any string input. Verified that compound scores for known positive phrases (e.g., `excellent product, highly recommend') are above 0.2 and for known negative phrases (e.g., `terrible quality, waste of money') are below -0.2.

  \item \textbf{Negation Handler:} Verified that a review containing `not bad' receives a lower compound score after negation handling than the raw VADER output, and that a review with no negation words receives the same score before and after the negation function.

  \item \textbf{Fake Review Detector (LinearSVC):} Verified that the \texttt{predict()} method returns either 0 (Genuine) or 1 (Fake) for any input string. Verified that a clearly generic computer-generated review is classified as Fake (CG) and a genuine product review is classified as Genuine (OR).

  \item \textbf{Aspect Sentiment Module:} Verified that reviews mentioning `battery backup is great' produce a positive aspect score for the Battery aspect. Verified that reviews mentioning `fast delivery' produce a positive score for the Delivery aspect.

  \item \textbf{Scraper Module:} Verified that the scraper returns a list of review strings for a known valid Flipkart product URL. Verified that Snap-path ChromeDriver binaries are correctly skipped and webdriver-manager is preferred.
\end{itemize}

\section{SECURITY TESTING}

The Opinion Miner platform accepts user-provided text and URLs, which introduces potential security considerations. The following aspects were tested to identify and mitigate vulnerabilities.

\subsubsection*{Security Aspects Tested:}

\begin{itemize}
  \item \textbf{Input Validation:} Verified that all user-submitted text inputs are processed as plain strings without executing embedded HTML tags or JavaScript content, preventing Cross-Site Scripting (XSS) vulnerabilities in the results display.

  \item \textbf{URL Validation:} Confirmed that URL inputs are validated for proper format and scheme (http/https) before being passed to the Selenium scraper, preventing arbitrary file access or server-side request forgery.

  \item \textbf{No Persistent Data Storage:} Confirmed that user-submitted text and scraped review data are not stored beyond the duration of the active Flask request session, ensuring user privacy.

  \item \textbf{Dependency Security:} All Python dependencies were reviewed for known CVEs. The \texttt{requirements.txt} specifies stable, widely-used library versions with no known active high-severity vulnerabilities.

  \item \textbf{Docker Security:} The Docker container runs as a non-root user where possible and does not expose any unnecessary ports beyond 5000.
\end{itemize}

\subsubsection*{Results:}

\begin{itemize}
  \item No critical security vulnerabilities were identified during testing.
  \item All user inputs are sanitized before processing and display.
  \item No sensitive data is stored or transmitted beyond the active session.
\end{itemize}

\section{INTEGRATION TESTING}

Integration testing was conducted after successful unit testing to verify the interactions between modules and ensure that the complete system operates as a coherent and reliable whole.

\subsubsection*{Objectives of Integration Testing:}

\begin{itemize}
  \item Validate proper communication and data format consistency between the Flask routes and each NLP module.
  \item Verify that the output of the Preprocessor module is in the correct format expected by the Sentiment Analysis module.
  \item Confirm that the FakeReviewDetector correctly passes only non-flagged reviews to the Sentiment pipeline and reports flagged reviews in the results summary.
  \item Verify that Aspect Analyser correctly receives genuine reviews and produces a structured per-aspect dictionary.
  \item Verify that chart images generated by the Visualisation module are correctly encoded as Base64 strings and successfully decoded and rendered in the frontend HTML response.
\end{itemize}

\subsubsection*{Integration Test Cases:}

\begin{itemize}
  \item Submit valid Flipkart URL $\rightarrow$ Scraper extracts reviews $\rightarrow$ FakeDetector filters reviews $\rightarrow$ Preprocessor cleans text $\rightarrow$ SentimentAnalyser scores reviews $\rightarrow$ AspectAnalyser produces 9-aspect breakdown $\rightarrow$ Aggregator computes summary $\rightarrow$ Results rendered in frontend. Full pipeline verified as functional.

  \item Enter negated review text (e.g., `not worth buying') $\rightarrow$ Preprocessor retains negation word $\rightarrow$ Sentiment Module applies dampening factor $\rightarrow$ Lower compound score returned than raw VADER output. Negation handling verified across pipeline.

  \item Enter review text classified as fake by LinearSVC $\rightarrow$ FakeDetector flags review $\rightarrow$ Sentiment pipeline not invoked for flagged review $\rightarrow$ Fake count incremented in results summary. Fake detection integration verified.

  \item Submit text input $\rightarrow$ Visualisation Module generates chart $\rightarrow$ Chart encoded as Base64 $\rightarrow$ Frontend renders chart inline in results section. Visualisation integration verified.
\end{itemize}

\subsubsection*{Integration Testing Results:}

\begin{itemize}
  \item Frontend-to-backend communication verified for all routes.
  \item API response times averaged below two seconds for text analysis and below ninety seconds for full URL analysis (50 pages).
  \item All module interfaces functioned consistently with expected data formats throughout the pipeline.
\end{itemize}

Testing at every stage confirmed that the platform is secure, reliable, and functionally complete. All major workflows were validated successfully, confirming readiness for deployment.

% ══════════════════════════════════════════════════════════════════════════════
%  CHAPTER 7: CONCLUSION AND FUTURE SCOPE
% ══════════════════════════════════════════════════════════════════════════════
\chapter{Conclusion and Future Scope}

\section{CONCLUSION}

Opinion Miner successfully demonstrates the development of a scalable, reliable, and interactive web-based system for the automated analysis of online product reviews. By integrating a real-time Selenium-powered scraping engine, a multi-stage text preprocessing pipeline, a VADER-based sentiment analysis model with custom negation handling, a Machine Learning-based fake review detection module (LinearSVC, 91.55\% accuracy), and a nine-aspect Aspect-Based Sentiment Analysis engine, the system provides a comprehensive end-to-end solution for opinion mining in the Indian e-commerce context.

The platform successfully addresses the core limitations of existing review analysis approaches. It eliminates the impracticality of manual review reading through automated scraping and processing. It goes beyond the reductionism of star-rating-based decisions by providing detailed sentiment scores, per-aspect breakdowns, review breakdowns, and visual distributions. It handles the pervasive problem of fake and deceptive reviews by filtering them using a proven ML classifier before sentiment computation, ensuring that only authentic customer feedback drives the final product verdict.

The Chrome Browser Extension brings the system's capabilities directly into the consumer's shopping workflow, enabling real-time sentiment analysis on any Flipkart product page without requiring a separate visit to the Opinion Miner web portal. The Docker containerisation and Render.com deployment ensure that the system is publicly accessible, maintainable, and reproducible across different environments without manual setup.

The system is built using industry-standard, widely adopted technologies including Python, Flask, Selenium, NLTK, scikit-learn, and Matplotlib, ensuring maintainability, extensibility, and ease of deployment. The modular architecture allows each component to be independently tested, updated, and scaled without disrupting the overall pipeline.

Comprehensive testing, encompassing system testing, unit testing, integration testing, and security testing, validated the functionality, accuracy, and robustness of all modules. As a result, Opinion Miner achieves its primary goal of providing a transparent, trustworthy, and practically useful opinion mining platform that bridges the gap between raw e-commerce review data and actionable consumer intelligence.

\section{FUTURE SCOPE}

While the current version of Opinion Miner fulfils its core objectives effectively, several opportunities for enhancement and expansion have been identified for future development:

\begin{itemize}
  \item \textbf{Expanded E-Commerce Platform Support:} The scraping and analysis pipeline can be extended to support additional Indian e-commerce platforms including Amazon India, Meesho, and Myntra, broadening the system's applicability beyond Flipkart.

  \item \textbf{Deep Learning Sentiment Analysis:} The VADER model, while efficient and practical for short informal text, can be supplemented or replaced with transformer-based models such as BERT (Bidirectional Encoder Representations from Transformers) or its domain-specific variants like BERTweet for higher accuracy on complex, sarcastic, or context-dependent review text.

  \item \textbf{Enhanced Multilingual Support:} Future work will involve integrating multilingual NLP models (XLM-RoBERTa or MuRIL) on a GPU-enabled server to enable full-fledged sentiment analysis for reviews written in Hindi, Telugu, Marathi, and Tamil, covering a significantly broader segment of Indian e-commerce consumers.

  \item \textbf{Product-to-Product Sentiment Comparison:} A comparison dashboard allowing users to input two or more Flipkart product URLs and view a side-by-side sentiment comparison across multiple product attributes would greatly enhance the platform's purchase decision support capabilities.

  \item \textbf{Advanced Analytics Dashboard:} A comprehensive analytics dashboard built with Streamlit or Dash can be developed to present sentiment evolution over time, top positive and negative keywords by product category, and comparative brand sentiment analysis using interactive time-series graphs and word clouds.

  \item \textbf{Mobile Application:} A cross-platform mobile application built using React Native or Flutter can be developed to bring Opinion Miner's review analysis capabilities to smartphone users, including an in-app browser view that triggers analysis when navigating to product pages.

  \item \textbf{Real-Time Monitoring and Alerts:} Scheduled monitoring of product sentiment with push notifications or email alerts when a product's sentiment drops below a configurable threshold would provide continuous consumer intelligence.

  \item \textbf{Automated Dataset Generation:} The user feedback collected through the interactive feedback loop can be used to automatically generate a labelled training dataset of correctly classified reviews, enabling periodic retraining of machine learning models to improve performance over time without manual annotation effort.

  \item \textbf{Marketplace API Integration:} Official Flipkart Affiliate API integration would replace Selenium-based scraping with a more stable, rate-limit-compliant data source, eliminating scraper fragility caused by UI changes.
\end{itemize}

The proposed enhancements will not only improve the accuracy, coverage, and user experience of the platform but will also make it significantly more comprehensive, scalable, and commercially viable, establishing Opinion Miner as a robust and trustworthy consumer intelligence tool for the Indian e-commerce ecosystem.

% ══════════════════════════════════════════════════════════════════════════════
%  REFERENCES
% ══════════════════════════════════════════════════════════════════════════════
\newpage
\begin{center}
{\large\textbf{References}}
\end{center}

\begin{singlespace}
\noindent [1] Hutto, C.J., \& Gilbert, E. (2014). VADER: A Parsimonious Rule-Based Model for Sentiment Analysis of Social Media Text. In \textit{Proceedings of the International AAAI Conference on Web and Social Media}, 8(1).

\medskip\noindent [2] Ott, M., Choi, Y., Cardie, C., \& Hancock, J.T. (2011). Finding Deceptive Opinion Spam by Any Stretch of the Imagination. In \textit{Proceedings of the 49th Annual Meeting of the Association for Computational Linguistics}.

\medskip\noindent [3] Mukherjee, A., Liu, B., \& Glance, N. (2012). Spotting Fake Reviewer Groups in Consumer Reviews. \textit{WWW Conference, Lyon, France}.

\medskip\noindent [4] Pang, B., \& Lee, L. (2008). Opinion Mining and Sentiment Analysis. \textit{Foundations and Trends in Information Retrieval}, 2(1-2), 1--135.

\medskip\noindent [5] Jindal, N., \& Liu, B. (2008). Opinion Spam and Analysis. \textit{ACM WSDM Conference}.

\medskip\noindent [6] Bird, S., Klein, E., \& Loper, E. (2009). \textit{Natural Language Processing with Python}. O'Reilly Media. Available at: https://www.nltk.org/

\medskip\noindent [7] Selenium Python Documentation. Selenium -- Browser Automation Framework. Available at: https://selenium-python.readthedocs.io/

\medskip\noindent [8] Flask Documentation. Flask -- A Lightweight Python Web Framework. Available at: https://flask.palletsprojects.com/

\medskip\noindent [9] Scikit-learn Documentation. Machine Learning in Python. Available at: https://scikit-learn.org/

\medskip\noindent [10] Matplotlib Documentation. Matplotlib -- Visualization with Python. Available at: https://matplotlib.org/

\medskip\noindent [11] Ravi, K., \& Ravi, V. (2015). A Survey on Opinion Mining and Sentiment Analysis: Tasks, Approaches and Applications. \textit{Knowledge-Based Systems}, 89, 14--46.

\medskip\noindent [12] Google. Chrome for Developers -- Manifest V3. Available at: https://developer.chrome.com/docs/extensions/mv3

\medskip\noindent [13] Docker Inc. Docker Documentation. Available at: https://docs.docker.com

\medskip\noindent [14] Vaswani, A., et al. (2017). Attention Is All You Need. \textit{Advances in Neural Information Processing Systems}.

\medskip\noindent [15] Pressman, R.S. (2019). \textit{Software Engineering: A Practitioner's Approach} (8th ed.). McGraw Hill Education.
\end{singlespace}

\end{document}



\documentclass[12pt,a4paper]{report}
\usepackage[top=1in, bottom=1in, left=1.5in, right=1.1in]{geometry}

\usepackage{times}
\usepackage{setspace}
\usepackage{titlesec}
\usepackage{tocloft}
\usepackage{graphicx}
\usepackage{amsmath}
\usepackage{booktabs}
\usepackage{array}
\usepackage{multirow}
\usepackage{fancyhdr}
\usepackage{etoolbox}
\usepackage{url}

%\usepackage[utf8]{inputenc}
%\usepackage[T1]{fontenc}
%\usepackage{fontspec}
%\usepackage{DejaVu Sans Mono}
%\setmonofont{DejaVu Sans Mono}
%\usepackage{verbatim}

\usepackage{fontspec}
\setmainfont{Times New Roman}
\setmonofont{DejaVu Sans Mono}


% Set line spacing
\onehalfspacing

% Configure page numbering
\pagestyle{fancy}
\fancyhf{}
\renewcommand{\headrulewidth}{0pt}

% For preliminary pages (roman numerals)
\fancypagestyle{preliminary}{
    \fancyhf{}
    \cfoot{\thepage}
    \renewcommand{\headrulewidth}{0pt}
}

% For main content (arabic numerals)
\fancypagestyle{main}{
    \fancyhf{}
    \cfoot{\thepage}
    \renewcommand{\headrulewidth}{0pt}
}

% Title formatting
\titleformat{\chapter}[display]
{\normalfont\huge\bfseries\centering}{\chaptertitlename\ \thechapter}{0.1pt}{\LARGE}

\titleformat{\section}
{\normalfont\bfseries\MakeUppercase}{\thesection}{1em}{}
\titleformat{\subsection}
{\normalfont\bfseries}{\thesubsection}{1em}{}

% Section spacing
\titlespacing*{\chapter}{0pt}{-30pt}{20pt}
\titlespacing*{\section}{0pt}{12pt}{6pt}
\titlespacing*{\subsection}{0pt}{12pt}{6pt}

% Paragraph formatting
\setlength{\parindent}{0pt}
\setlength{\parskip}{12pt}

% === Border on every page (safe version) ===
\usepackage{tikz}
\usepackage{eso-pic}

\makeatletter
\AddToShipoutPictureBG{%
  \AtPageLowerLeft{%
    \begin{tikzpicture}[remember picture, overlay]
      % Left margin = 1.5in, right/top/bottom = 1in
      \draw[line width=1pt, color=black]
        ([xshift=1.3in,yshift=0.9in]current page.south west)
        rectangle
        ([xshift=-0.9in,yshift=-0.9in]current page.north east);
    \end{tikzpicture}%
  }
}
\makeatother

